\documentclass{./EasyChair3.5/easychair}
\usepackage{unicode-math}
\usepackage{multicol}
\usepackage{mathtools}
\usepackage{csquotes}
\usepackage{tikz}
\usetikzlibrary{backgrounds,calc,positioning,arrows.meta,quotes,cd}
\usepackage{amsthm}
\usepackage{caption}
\usepackage[inline]{enumitem}
\usepackage{hyperref}

\newtheorem{theorem}{Theorem}
\theoremstyle{definition}
\newtheorem{definition}{Definition}
\newtheorem{assumption}{Assumption}
\newtheorem{convention}{Convention}
\newtheorem{observation}{Observation}
\newtheorem{corollary}{Corollary}
\theoremstyle{remark}
\newtheorem{notation}{Notation}
\newtheorem{remark}{Remark}
\usepackage{microtype}
\usepackage{newunicodechar}
\usepackage{mathtools}
\newunicodechar{⇒}{\ensuremath{\Rightarrow}}
%\DeclareUnicodeCharacter{21D2}{\ensuremath{\Rightarrow}}
\newunicodechar{ℰ}{\ensuremath{\mathscr{E}}}
\newunicodechar{₀}{\ensuremath{{}_{0}}}
\newunicodechar{▷}{\ensuremath{\triangleright}}
\newunicodechar{ℕ}{\ensuremath{ℕ}}
\newunicodechar{↦}{\ensuremath{↦}}
\newunicodechar{∘}{\ensuremath{∘}}
\newunicodechar{ᵢ}{\ensuremath{_i}}
\newunicodechar{≔}{\ensuremath{≔}}
\newunicodechar{≈}{\ensuremath{\approx}}
\newunicodechar{₁}{\ensuremath{{}_{1}}}
\newunicodechar{₂}{\ensuremath{{}_{2}}}
\newunicodechar{λ}{\ensuremath{\lambda}}
\newunicodechar{≺}{\ensuremath{<}}
\newunicodechar{∈}{\ensuremath{\in}}
\newunicodechar{⊥}{\ensuremath{\bot}}
\newunicodechar{ε}{\ensuremath{\varepsilon}}
\newunicodechar{↓}{\ensuremath{\downarrow}}
%% Invisible characters for otherwise clashing definitions.
\newunicodechar{⇃}{\ensuremath{\!}}
\newunicodechar{⇂}{\ensuremath{\!}}
\newunicodechar{↿}{\ensuremath{\!}}
\newcommand{\sle}{\sqsubseteq} %sort \le
\newcommand{\sgt}{\sqsupset} %sort \ge
\newunicodechar{⊒}{\ensuremath{\sgt}}
\newunicodechar{⊑}{\ensuremath{\sle}}
\newunicodechar{ₗ}{\ensuremath{_l}}
\newunicodechar{ᵣ}{\ensuremath{_r}}
\newunicodechar{⁻}{\ensuremath{^-}}
\newunicodechar{¹}{\ensuremath{^1}}
\newunicodechar{∷}{\ensuremath{\mathbin{::}}}
\newunicodechar{┌}{|\ensuremath{\lceil}}
\newunicodechar{┐}{\ensuremath{\rceil}|}
\newunicodechar{₁}{\ensuremath{_1}}
\newunicodechar{₂}{\ensuremath{_2}}
\newunicodechar{→}{\ensuremath{\to}}
\newunicodechar{⸴}{\ensuremath{,}}
\newunicodechar{ℬ}{\ensuremath{\mathcal{M}\!}}
\newunicodechar{⊔}{\ensuremath{\sqcup}}
\newunicodechar{⊓}{\ensuremath{\sqcap}}
\newunicodechar{⊎}{\ensuremath{\uplus}}
% elide "on obj/on morphs" distinction in code presentation
\newunicodechar{₀}{}
\usepackage{../agda}
\usepackage{catchfilebetweentags}

\DeclarePairedDelimiter{\expi}{\lceil}{\rceil}

\newcommand{\nicehref}[2]{%
\begin{tikzpicture}[baseline=(txt.base)]
\node[text=blue!80!white!40!black,inner sep=0pt,text depth=1.1pt] (txt)
  {\href{#1}{{#2}{\hspace*{1pt}\raisebox{.0pt}{\color{blue!30!white}\tiny\faExternalLink}}}};
% the text depth (space below the baseline) determines the position of the (following) underline.
% I've tried to align it with underscore characters.
\begin{scope}[on background layer]
\draw[color=blue!25!white] (txt.south west) -- (txt.south east);
\end{scope}
\end{tikzpicture}%
}

\newcommand{\inj}{\mathsf{in}}
\newcommand{\C}{{\ensuremath{\mathcal{C}}}}

\title{Intrinsically correct \& total functional programming with enrichment, contractive functors, recursive coalgebras and cofree presheaves}

\author{
  Cass Alexandru \inst{1} \and Steven Schaefer \inst{2} \and Max S. New \inst{2}
}

\institute{
  RPTU Kaiserslautern-Landau, Radboud University Nijmegen
  \and University of Michigan
}

\authorrunning{Alexandru, Schaefer, New}

% \titlerunning{} has to be set to either the main title or its shorter
% version for the running heads. When processed by
% EasyChair, this command is mandatory: a document without \titlerunning
% will be rejected by EasyChair
\titlerunning{Intrinsically correct \& total fp}
\begin{document}

\maketitle

Recent progress has been made \cite{alexandruIntrinsicallyCorrectAlgorithms2026} in formalising synergies between proofs of partial and total correctness of algorithms defined using structured recursion \cite{hinzeConjugateHylomorphismsMother2015}.
This was acheived with the introduction of the novel notion of a well founded functor.
We acheive the same result, i.e. a criterion for all coalgebras for an endofunctor on the power category \(𝒞^A\), for \(A\) the carrier of a well founded ordering \(<\), to be recursive.
Our criterion is a novel instantiation of the preexisting notion of a locally contractive functor \cite{birkedalFirstStepsSynthetic2012}.
This instantiation is acheived by enriching \(𝒞^A\) in presheaves over \((A,<)\) using change of base of enrichment along the cofree presheaf functor.
Our hope is that demonstrating how this preexisting notion can be instantiated for this novel application may lead to a broader adoption of the technique, as well as enabling reuse of all existing known results about locally contractive functors in this setting.

\section{Background}
\label{sec:background}

\paragraph{Recursive Coalgebras}
Expressing recursive algorithms using coalgebra-to-algebra morphisms has a long history \cite{CaprettaUV06,jeanninWellfoundedCoalgebrasRevisited2017,taylor99,meijerFunctionalProgrammingBananas1991}.
The key idea is the following: In a base category \(𝒞\) modeling the types and functions of your functional programming language, choose an endofunctor \(F \colon 𝒞 → 𝒞\). We can think of this functor as describing a distribution; e.g.\ the functor \(B X ≔ 1 + X × A × X\) describes the distribution of either a base case \(* : 1\) or two inputs \(l , r : X\) and an output \(o : A\).
A recursive algorithm \(h\) can be seen as the following sequence of steps:
\begin{enumerate*}[itemjoin=~]
\item \emph{divide} input into new input, with a distribution given by \(F_0\);
\item recur under \(F_1\);
\item \emph{combine} the results of the recursive invocations.
\end{enumerate*}
E.g.\ the recursive case of Quicksort \autoref{fig:qsort-recursive} (due to \cite{CaprettaUV06}); the general case is \autoref{fig:coalg-alg-morphism}.

\noindent
\begin{minipage}[t]{.6\textwidth}
\begin{tikzcd}[ampersand replacement=\&,
every arrow/.append style={mapsto},column sep=2cm]
 {(i_{≤ p}, p, i_{> p})} \arrow[r, "\mathsf{qsort} \times \mathrm{id}\times \mathsf{qsort}"] \& {(\mathsf{qsort}(i_{≤ p}), p, \mathsf{qsort}(i_{> p}))} \arrow[d, "\mathsf{concat}"] \\
ip \arrow[u,"\mathsf{partition}"] \arrow[r,"\mathsf{qsort}"] \&  {\mathsf{qsort}(i_{≤ p})\, p\, \mathsf{qsort}(i_{> p})}
\end{tikzcd}\vspace{-1\baselineskip{}}
\captionof{figure}{quicksort recursive case}\label{fig:qsort-recursive}
\end{minipage}\hfill
\begin{minipage}[t]{.4\textwidth}
\begin{tikzcd}
F_0I \arrow[r, "F_1h"] & F_0O \arrow[d, "a"] \\
I \arrow[u,"c"] \arrow[r,"h"] &  O
\end{tikzcd}
\makebox[\linewidth][r]{\parbox{1.25\linewidth}{\captionof{figure}{coalgebra-to-algebra morphism}\label{fig:coalg-alg-morphism}}}
\end{minipage}

A coalgebra \(c\) is called \emph{recursive}, if for every target algebra \(a\) there is a unique \(h\) such that \eqref{fig:coalg-alg-morphism} commutes, i.e.\ the recursive equation \(h = c ; F h ; a\) has a unique solution and can be read as a \emph{definition}.
Intuitively, this means \(c\) is a division step that, when iterated, builds only well founded trees of shape \(F\), regardless of the starting input \(i : I\).
The challenge of defining an algorithm as a coalgebra-to-algebra morphism in a total language then can be reduced to proving the recursivity of its coalgebra.

This challenge becomes easier to tackle when we take partial correctness into account. Continuing with our Quicksort, we home in on its coalgebra, \textsf{partition}. Its type, \(\mathsf{List} A → 1 + \mathsf{List} A × \mathsf{List} A\) is really too coarse. The partial correctness of \textsf{partition} entails
\begin{enumerate*}[itemjoin=~]
\item the output \(B (\mathsf{List} A)\) value contain the same \emph{elements} as the input;
\item the elements to the left/right of the selected pivot be respectively smaller/greater than it.
\end{enumerate*}
One approach to refine its type in order to capture this, due to \cite{AlexandruCRW25}, is to move to the setting of \emph{families} in \textsf{Set} indexed by the set \(ℳA\) of finite multisets of element type \(A\).
The refined \(B\) becomes \(\bar{B}\) (expressed for better comprehensibility in the slice category \(\mathsf{Set}/ℳA\) instead of in \(\mathsf{Set}^{ℳA}\)) :
\[
  \bar{B} \begin{pmatrix}
    X \\
    f
  \end{pmatrix}
  ≔
  \begin{pmatrix}
    1 \\
    * ↦ \emptyset
  \end{pmatrix}
  +
  \begin{pmatrix}
    \{(l,p,r) ∈ X × A × X \mid ∀ e_l ∈ f(l), e_r ∈ f(r).\, e_l ≤ p ≤ e_r\} \\
    (l, p, r) ↦ f(l) \uplus \{p\} \uplus f(r)
  \end{pmatrix}
\]
The interesting thing to note is the following: The multiset indices in functorial positions are smaller than the result index: \(f(l),f(r) < f(l) \uplus \{p\} \uplus f(r)\).
In particular, that means we are able to define a functorial action of \(F\), \(F_1'\) that, intuitively, takes morphisms \(A → B\) at smaller indices and produces a morphism \(F_0 A → F_0 B\) at the current index, such that \(F_1'\) agrees with \(F_1\). This then allows us to essentially construct a solution to \eqref{fig:coalg-alg-morphism} as an inhabitant of the homset \(I → O\) using well founded recursion. We proceed to make all of these notions formal. For notions of enriched category theory we mostly follow~\cite{weideUnivalentEnrichedCategories2026}.

\paragraph{Enrichment}
To talk about an action of a functor \(F\) that isn't \(F_1\), we need the notion of:
\begin{definition}[Enrichment of a category \(𝒞\) in a monoidal category \(𝒱\), {\cite[Def.~2.2]{weideUnivalentEnrichedCategories2026}}]
  An enrichment of a category \(𝒞\) in a monoidal category \(𝒱,⊗,I\) consists of the following data:
  \begin{itemize}
  \item for every \(x, y \colon 𝒞\) the external Hom \(ℰ_{X,Y}\), an external identity morphism, sequencing operation, and external counterparts to unitality \& associativity coherences
  \item for every \(x, y \colon 𝒞\) a bijection \(\expi{-}\colon 𝒞(x,y) ≅ 𝒱[I,ℰ_{X,Y}]\) witnessing that \(𝒞\) and the \emph{enriched category} with objects \(𝒞_0\) and homsets \(𝒱[I,ℰ_{X,Y}]\) are isomorphic as categories.
  \end{itemize}
\end{definition}
\begin{definition}[Enrichment of a functor, {\cite[Def.~2.4]{weideUnivalentEnrichedCategories2026}}]
  Given categories \(C, D\) and enrichments \(EC, ED\) in a monoidal category \(V\) an \emph{enrichment} of a functor \(F \colon C → D\) is an external functorial action
  \(F_1' \colon 𝒱[EC.ℰ_{X,Y},ED.ℰ_{F_0 X, F_0 Y}]\), obeying external counterparts to preservation of identities and sequencing, such that \(F_1\) and \(F_1'\) \emph{agree}:
  \[∀ x, y : C_0,\, f : C[x,y].\ ED.\expi{F_1 f} = EC.\expi{f} ; F_1' \]
\end{definition}

\paragraph{Locally Contractive Functors}
\begin{definition}[{\cite[Def.~7.2]{birkedalFirstStepsSynthetic2012}}]
  Let \(V\) be a model of guarded recursive terms \cite[Def.~6.1]{birkedalFirstStepsSynthetic2012}. Given categories \(C,D\), with enrichments \(EC, ED\) in \(V\), a functor \(F \colon C → D\) is \emph{locally contractive} if it factors as the sequence of the enrichments
  \[CE \xrightarrow{\mathsf{next}} ▷CE \xrightarrow{F'} ED\]
\end{definition}

\begin{corollary}
  Let \(V\) be a model of guarded recursive terms, \(C,D\) categories with enrichments \(EC, ED\) in \(V\). A functor is \emph{locally contractive} if it as an enrichment of the sequence of enrichments \(CE \xrightarrow{\mathsf{next}} ▷CE \xrightarrow{F'} ED\), i.e.
\[∀ x, y : C_0,\, f : C[x,y].\ ED.\expi{F_1 f} = EC.\expi{f} ; \mathsf{next} ; F_1' \]
\end{corollary}


\paragraph{Change of Base}

\paragraph{Example Application to QuickSort}

\bibliographystyle{plain}
\bibliography{../works}

\end{document}
